\documentclass[a4paper]{article}

\usepackage[fontset=fandol]{ctex}
\usepackage{amsmath, amssymb}
\usepackage{graphicx}
\usepackage[margin=2.5cm]{geometry}
\usepackage[most]{tcolorbox}
\usepackage{etoolbox}
\usepackage{listings}
\usepackage{hyperref}
\usepackage{booktabs}
\usepackage{subcaption}
\usepackage{float}
\usepackage{tikz}
\IfFileExists{pgfplots.sty}{
    \usepackage{pgfplots}
    \pgfplotsset{compat=1.18}
}{}

\graphicspath{{figures/}}

\newtcolorbox{knowledgebox}[1]{
    enhanced, colback=blue!5!white, colframe=blue!75!black, colbacktitle=blue!75!black,
    coltitle=white, fonttitle=\bfseries, title=#1, attach boxed title to top left={yshift=-2mm, xshift=2mm},
    boxrule=1pt, sharp corners
}

\newtcolorbox{importantbox}[1]{
    enhanced, colback=yellow!10!white, colframe=yellow!80!black, colbacktitle=yellow!80!black,
    coltitle=black, fonttitle=\bfseries, title=#1, sharp corners
}

\newtcolorbox{warningbox}[1]{
    enhanced, colback=red!5!white, colframe=red!75!black, colbacktitle=red!75!black,
    coltitle=white, fonttitle=\bfseries, title=#1, sharp corners
}

\newcommand{\notetitle}{MIT RES.9-009 第02讲：Neuroblox 图形界面}
\newcommand{\noteauthors}{MIT RES.9-009 课程整理 \& Codex}
\newcommand{\notedate}{\today}
\newcommand{\videochannel}{MIT Video Productions；讲者 Helmut H. Strey}
\newcommand{\videopublishdate}{2025-04-10（讲课日期：2025-01-13）}
\newcommand{\videoduration}{20:51}
\newcommand{\videourl}{https://www.youtube.com/watch?v=XlBRJps84zE}
\newcommand{\videocoverpath}{figures/cover.png}

\begin{document}

\begin{titlepage}
\centering
{\Large 课程笔记\par}
\vspace{1.2cm}
{\huge\bfseries \notetitle\par}
\vspace{0.8cm}
{\large \noteauthors\par}
\vspace{0.3cm}
{\large \notedate\par}
\vspace{1.2cm}

\ifdefempty{\videocoverpath}{
{\small 请在生成笔记时填入视频封面图的本地路径。\par}
}{
\includegraphics[width=0.82\textwidth,height=0.45\textheight,keepaspectratio]{\videocoverpath}\par
}

\vfill
\begin{tcolorbox}[width=0.9\textwidth, colback=black!2!white, colframe=black!60, sharp corners]
\textbf{视频作者/频道}：\videochannel\par
\textbf{发布日期}：\videopublishdate\par
\textbf{视频时长}：\videoduration\par
\textbf{视频链接}：
\ifdefempty{\videourl}{
未填写
}{
\href{\videourl}{\nolinkurl{\videourl}}
}
\end{tcolorbox}
\end{titlepage}

\tableofcontents
\newpage

\begin{importantbox}{证据层说明}
本讲没有找到独立的官方 slide PDF。正文主要依据字幕、官方课程页和视频帧；课程页中的 logo 首图只承担品牌展示，不承担本讲教学内容，因此已在 \texttt{omission\_log.jsonl} 中明确记录。
\end{importantbox}

\section{GUI 入口、flat brain 与协作状态}
这一讲一开始就把两件事说得很明确。第一，演示的是 Neuroblox 的新 GUI，而不是 JuliaHub 上那个完全稳定的版本，因为还有一些功能没有完全接通。第二，Neuroblox 现在已经进入可以给学生试用的阶段，但它仍然是一个快速迭代的系统，所以讲者反复提醒，遇到 bug 不要惊讶，要把它当成真实协作的一部分。

官方课程页在这个层面上也给了同样的背景：Neuroblox 被定位成一个基于 Julia 的模块化平台，支持从单个 block 到更大系统的拼接，并允许通过 graphical interface 直接组装模型。第 2 讲的作用，就是把这个抽象定位落到界面上，告诉你到底怎么点、怎么连、怎么跑。

\begin{figure}[H]
\centering
\includegraphics[width=\textwidth]{gui_overview.png}
\caption{GUI 总览：flat brain、节点区和右侧属性面板同时可见。这个画面对应讲者解释 canvas、sidebar 和 simulation 面板的段落（约 00:08:05）。}
\label{fig:gui_overview}
\end{figure}

最核心的交互语义其实很简单，但必须记牢。canvas 上的脑区是一个 flat map，也就是把 brain region 展开到平面上；把 block 放到某个区域上以后，你就可以通过 Alt 或 Option 建边。左键点击是选择，左键拖拽边上的 handle 可以改变边的几何形状。换句话说，GUI 里既有“结构放置”，也有“连接编辑”，这两者不是同一个动作。

\begin{knowledgebox}{几个必须记住的操作语义}
Alt 或 Option 用来建边；左键点 node 是选 node；左键点 edge 是选 edge；拖拽 edge anchor 可以改边形状。这个交互设计的目的，是让你在不离开 flat brain 的情况下同时编辑结构和动力学。
\end{knowledgebox}

右侧面板的结构也很明确。\texttt{Nodes \& Edges} 用来改节点和连接的属性，\texttt{Params} 用来集中管理参数，\texttt{Sim} 负责仿真设置，\texttt{Parameter Fitting} 和 \texttt{Reinforcement Learning} 目前都还不够稳定，\texttt{Data Loader} 则负责加载数据。讲者特地提醒，不要把还没完全工作的功能当成课程主流程的一部分。

\begin{warningbox}{当前可用性提醒}
讲者明确说过，\texttt{parameter fitting} 目前不建议依赖，\texttt{reinforcement learning} 也还没有在这个 GUI 上形成稳定经验。教学上应把它们当成未来功能，而不是当前练习的主线。
\end{warningbox}

另一个很现实的细节是云端启动时间。连接 JuliaHub 后，系统要在 AWS 里预留计算资源，随后 Julia 又要进行 JIT 编译，所以第一次点开和第一次跑仿真会慢很多。讲者把这一点解释得很直白：不是 GUI “坏了”，而是云服务和编译开销在前面排队。

你还可以把这一讲里的平台架构理解成三层协作：前端负责界面，graphics server 负责作图，Julia backend server 负责真正的计算。讲者说得很清楚，这三者是互相通信的；等仿真结束，数据先传给图形服务器，再回到界面里显示。这也是为什么 GUI 看起来是“点一下就出图”，但背后其实是一条很长的异步链路。

\subsection{本章小结}
这一部分不是在讲“GUI 好不好看”，而是在讲 Neuroblox 的操作语义：flat brain 负责空间布局，Alt/Option 负责连边，右侧面板负责参数与仿真，云端和 Julia 则负责执行。把这些语义记牢，后面的示例才有操作基础。

\section{Harmonic Oscillator：最小可运行的双节点系统}
讲者接下来选了一个最小但很适合练手的例子：harmonic oscillator neural mass model。它足够简单，能让你快速摸到 GUI 的所有关键操作；它也足够接近动力系统的本体，让你理解为什么 Neuroblox 不是“画图工具”，而是“搭动力学系统的工具”。

最基本的连续形式可以写成
\[
\dot{x}(t)=y(t), \qquad \dot{y}(t)=-\omega^2 x(t)-\gamma y(t)+I_{\mathrm{couple}}(t).
\]
如果只看无阻尼版本，那么把一阶系统消元后就是
\[
\ddot{x}(t)+\omega^2 x(t)=0,
\]
解是正弦或余弦，频率由 $\omega$ 控制。讲者进一步提醒，若只有阻尼项而没有合适的耦合，这个系统会衰减掉，不会形成你想要的稳定振荡。所以在 GUI 里，他不是只放一个 oscillator，而是放两个互相作用的 oscillator。

\begin{figure}[H]
\centering
\includegraphics[width=\textwidth]{harmonic_pair.png}
\caption{两个 harmonic oscillator，HO1 和 HO2 已经放到 canvas 上，并且能在右侧看到边的属性面板。这个状态对应讲者说明 self-interaction、Alt 拖拽和命名的片段（约 00:10:02）。}
\label{fig:harmonic_pair}
\end{figure}

讲者给出的实操顺序很稳定：先在底部 block 面板里找到 \texttt{HarmonicOsc}，拖到 canvas 上，重命名为 \texttt{HO1} 和 \texttt{HO2}，然后给每个 node 加 self-interaction，再通过 Alt/Option 从一个 node 连到另一个 node。起始参数并不是理论定值，而是一个经验起点。讲者建议先把 self-interaction 设成 inhibitory，也就是大约 \texttt{-1}，再把跨节点连接设成更强一些的正值，例如 \texttt{5} 和 \texttt{6}，然后观察是否能把系统推到振荡区间。

\begin{importantbox}{为什么要用两个 oscillator}
单个 damped oscillator 更像一个会衰减的振动器；两个相互作用的 oscillator 才更像一个可以通过耦合维持节律的系统。讲者在这里的重点不是某个固定参数，而是让你感受“自耦合 + 互耦合”如何改变系统的相图。
\end{importantbox}

GUI 里还给了一个很有用的事实：Neuroblox 的时间单位是 milliseconds，所以输入的 \texttt{1} 不是秒，而是 \texttt{1 ms}。仿真面板还允许你设定 duration、relative tolerance、absolute tolerance，以及 stiff / non-stiff solver。讲者的建议是，先让系统跑通，再去调这些数值细节；第一次运行还会慢一些，因为 Julia 需要先编译相关包，而云端环境里这一步往往最耗时间。

当仿真跑完以后，右侧的结果面板会显示轨迹，你可以把图放大、挪到另一侧、切换查看 node 的参数，甚至打开 \texttt{Show latex} 看系统方程。讲者还顺手提到，一个想当然的小问题其实很适合作为练习：如果你已经把这个系统调到会振荡，如何把它的频率加倍？这个问题没有在现场给出唯一答案，它本来就是为了让你理解参数、耦合和频率之间的关系。

\subsection{本章小结}
这一部分的价值在于把 GUI 操作和动力系统直觉绑在一起。你不是在“摆两个图标”，而是在构造一个能通过 self-interaction 和 cross-coupling 进入振荡区间的最小系统。会搭这个例子，后面的 cortical demo 才不是黑箱。

\section{Cortical demo：从单块到双区域}
第三段把系统规模往上推了一层。讲者先新建一个 \texttt{cortical demo} 项目，然后从 composite blocks 里选出 \texttt{Cortical} block，把它放到 canvas 上并命名为 \texttt{C1}。这个例子说明了一件很重要的事：Neuroblox 的核心不是单个 neuron，而是 \texttt{blocks of blocks}。你可以把 WTA、cortical region、stimulus block 继续封装，再把这些封装好的块放进更大的系统。

在参数层面，\texttt{Cortical} block 会出现若干控制项。最常见的有 \texttt{N\_wta}、\texttt{N\_exci}、\texttt{G\_syn\_exci}、\texttt{G\_syn\_inhib} 和 \texttt{G\_syn\_ff\_inhib}。讲者的实用建议很直接：如果默认的 \texttt{N\_wta = 10} 太慢，就先把它减到 \texttt{5}，因为你当前的目标是理解工作流，不是让一次演示把时间烧完。

\begin{figure}[H]
\centering
\includegraphics[width=\textwidth]{cortical_demo.png}
\caption{双 cortical region 的结果视图：左侧显示 adjacency matrix / cross correlation 类输出，右侧是 \texttt{C2} 的参数面板。这个画面对应 00:19:22 左右的演示。}
\label{fig:cortical_demo}
\end{figure}

仿真完成后，GUI 提供的不只是一个时间序列，还包括多个特殊视图：\texttt{raster plot}、\texttt{stack plot}、\texttt{adjacency matrix}、\texttt{power spectrum}。讲者强调，\texttt{adjacency matrix} 里看到的连接是按 density 随机采样出来的，density 控制的是连接概率，不是突触权重；权重仍然要另行设置。这个区别非常重要，因为它直接决定你是在调结构稀疏性，还是调动力学强度。

\begin{knowledgebox}{density 不是 weight}
\texttt{density} 决定“有多少条边会被采样出来”，\texttt{weight} 决定“已经采样出来的边有多强”。这两个概念必须分开，不然你会把结构问题和动力学问题混成一件事。
\end{knowledgebox}

讲者随后把系统扩展成两个 cortical regions，\texttt{C1} 和 \texttt{C2}。在这个过程中，讲者又一次强调了 GUI 的实用细节：如果结果图还停留在旧参数上，就再点一次 \texttt{Run simulation}；如果 region 太大，仿真会明显变慢；如果想看更深的结构，还可以打开 \texttt{Show latex} 查看方程。更进一步，如果你愿意，Neuroblox 还可以把 image stimulus block 接到 visual cortex，再把信号送到 prefrontal cortex。讲者把这部分留给后面愿意继续试的人，而不是把它当成今天必须完成的任务。

\subsection{本章小结}
这一段把 GUI 从“单个 oscillator 的练习场”推进成了“可拼装的 cortical system”。\texttt{Cortical} block、\texttt{C1/C2}、\texttt{adjacency matrix} 和 \texttt{power spectrum} 共同说明：Neuroblox 的可视化不是装饰，而是模型结构本身的一部分。

\section{总结与延伸}
这一讲真正教你的不是某个按钮，而是一整套操作顺序。先在 flat brain 上放 block，再用 Alt/Option 建边；先用 harmonic oscillator 验证你是否理解 node / edge / parameter 的关系，再把同样的逻辑推广到 cortical region 和 composite blocks。也就是说，GUI 不是课外附带工具，而是 Neuroblox 的建模入口。

如果把这讲压成三句话，就是：
\begin{itemize}
  \item \texttt{flat brain} 让你把脑区、block 和连接放到同一个平面上编辑。
  \item \texttt{HarmonicOsc} 让你验证自己真的会建边、调参、跑仿真。
  \item \texttt{Cortical} block 让你看到单个 block 如何被组合成更大的回路系统。
\end{itemize}

这也是为什么讲者最后反复鼓励大家“多玩玩 GUI”。对后面的教程来说，能不能顺利进入 \texttt{WinnerTakeAll}、\texttt{Cortical}、\texttt{ImageStimulus} 等示例，取决于你是否已经把这一讲里的交互语义掌握住了。

\begin{importantbox}{课程位置}
官方课程总览里，\texttt{Blox and Connections in Neuroblox} 之后紧接着就是 \texttt{Neurons, Neural Masses and Sources}、\texttt{Circuit Models}、\texttt{Decision Making} 等单元。本讲的作用，正是把你从“知道有这些单元”推进到“知道怎么真的在 GUI 里搭起来”。
\end{importantbox}

\subsection{本章小结}
第 2 讲是后续所有教程的操作底座。它把界面、参数、仿真和复合块串成了一条完整的手工流程；只要这条链路打通，后面的更大回路和更复杂实验就只是规模上的扩展，而不是概念上的断层。

\end{document}
